\subsection{扩张的次数}

设 A 是域 K 上的代数。它特别是 K 上的向量空间；这个向量空间的维数称为 A 在 K 上的次数，记作 \([A:K]\) （II，第 293 页）。根据定义， \([A:K]\) 因此是 A 在 K 上的任意基的基数。这个定义特别适用于 K 的扩张的情形。

次数为 1 的扩张是平凡的。次数为 2（分别为 3 等）的扩张称为二次（分别为三次等）扩张。有限次数的扩张有时被滥用语言地称为有限扩张。

定理 1. --- 设 E 是 K 的扩张，A 是 E 上的代数。那么我们有 \([A:K]=[A:E]\)。\([E:K]\)。特别地，若 F 是 E 的扩张，我们有

\[
[ \mathbf {F}: \mathbf {K} ] = [ \mathbf {F}: \mathbf {E} ]. [ \mathbf {E}: \mathbf {K} ].\tag{1}
\]

该定理只是 II，第 222 页，命题 25 的一个特例；更确切地说，如果 \((a_{\lambda})_{\lambda \in L}\) 是 A 在 E 上的一组基，且 \((b_{\mu})_{\mu \in M}\) 是 E 在 K 上的一组基，则族 \((a_{\lambda}b_{\mu})_{(\lambda ,\mu)\in L\times M}\) 是 A 在 K 上的一组基。

推论 1. --- 设 K、E、F 为三个域，使得 \(K \subset E \subset F\) 且 \([F:K]\) 是有限的。则次数 \([E:K]\) 与 \([F:E]\) 是 \([F:K]\) 的因子。

如果次数 \([F:K]\) 是素数，则 F 除 K 和 F 外没有其他子扩张。但要注意，当 \([F:K]\) 不是素数时，不一定存在 F 的除 K 和 F 以外的子扩张（参见 V，第 146 页，习题 1）。

推论 2. --- 设 K、E 和 F 为三个域，使得 \(K \subset E \subset F\) 。假设 \([F:K]\) 是有限的，则关系 \([E:K] = [F:K]\) 等价于 E = F，且关系 \([F:E] = [F:K]\) 等价于 E = K。

因为若 L 是 \(L'\) 的一个扩张域，则 \([L:L'] = 1\) 等价于 \(L' = L\) 。

命题 1. --- 设 A 是域 K 上的一个有限次代数。若一个元素 \(a \in \mathbf{A}\) 在 A 中不是左（相应地，右）零因子，则它在 A 中可逆。

因为 A 在 K 上的向量空间按假设是有限维的，且线性映射 \(x \mapsto ax\)（相应地 \(x \mapsto xa\)）（从 A 到 A）是单射；因此它是双射（II，第 298 页，推论），从而（I，第 16 页，注记）a 在 A 中可逆。

推论。--- 设 A 是域 K 上的一个有限次交换代数。若 A 是整环，则它是一个域。

